\chapter{Limping in Children}

\section*{Definition}
Limping is an altered walking pattern caused by pain, weakness, restricted movement, injury, or neurologic dysfunction. A new limp without an obvious minor injury deserves assessment.

\section*{When to See a Doctor}
Emergency evaluation is needed for inability to walk, severe pain, fever, a swollen joint, deformity, night pain, weight loss, or a child who appears acutely ill. Arrange prompt evaluation for persistent, recurrent, unexplained, or function-limiting symptoms.

\section*{How It Develops}
The symptom results from irritation, inflammation, altered organ function, injury, infection, or disruption of normal neurologic or vascular signaling. Age, pregnancy status, onset, distribution, severity, and associated findings determine the urgency and likely causes.

\section*{Causes}
\begin{itemize}
  \item Minor trauma, transient synovitis, infection, fracture, juvenile arthritis, Perthes disease, slipped capital femoral epiphysis, developmental hip disease, and neurologic or muscular disorders.
  \item Medication effects, environmental exposures, and underlying chronic disease may also contribute.
  \item Serious causes are less common but must be considered when symptoms are sudden, progressive, severe, or associated with systemic illness.
\end{itemize}

\section*{Related Symptoms}
\begin{itemize}
  \item Pain, fever, fatigue, nausea, weakness, or changes in normal function
  \item Bleeding, swelling, discharge, breathing difficulty, or neurologic symptoms when a serious cause is present
\end{itemize}

\section*{Medical Evaluation}
\begin{itemize}
  \item History addresses onset, duration, severity, triggers, injuries, exposures, medications, medical conditions, age, and pregnancy status when relevant.
  \item Examination includes vital signs and focused assessment of the relevant organ system, neurologic status, hydration, and function.
  \item Testing may include blood or urine studies, cultures, imaging, physiologic testing, fetal or pediatric assessment, or specialist examination.
\end{itemize}

\section*{Treatment Options}
Treatment is directed at the cause and may include supportive care, targeted medication, treatment of infection or inflammation, correction of metabolic disease, physical therapy, or specialist procedures. Persistent symptoms require reassessment.

\section*{Self-Care and Home Management}
Avoid known triggers, protect the affected area, maintain appropriate hydration, and record timing and associated symptoms. Use only age- or pregnancy-appropriate medicines recommended by a clinician. Do not delay emergency care while trying home treatment.

\section*{References}
\begin{thebibliography}{99}
\bibitem{limping_in_children:limping} Bono KT, Squires J, Della Morte M. Evaluation of the limping child. \textit{Am Fam Physician}. 2015;91:611--618.
\bibitem{limping_in_children:kocher} Kocher MS, Zurakowski D, Kasser JR. Differentiating between septic arthritis and transient synovitis of the hip in children: an evidence-based clinical prediction algorithm. \textit{J Bone Joint Surg Am}. 1999;81:1662--1670.
\bibitem{limping_in_children:leet} Leet AI, Skaggs DL. Evaluation of the acutely limping child. \textit{Am Fam Physician}. 2000;61:1011--1018.
\end{thebibliography}
